\section{Discussion}

Under the subject-disjoint protocol the Mamba classifier is the best model we
tested, but the margin over a well-tuned Random Forest is modest: five accuracy
points on both feet, and the two are level on AUC. The larger lesson of these
experiments is not that a state-space model wins. It is that two things we
initially treated as details -- how the signal is presented to the network, and
how the data is split -- moved the numbers far more than the choice of
architecture did.

The failure with raw input was what made the first point concrete. On folds~1
and~4 training settled at a validation AUC of about $0.53$, while the same code
on folds~2, 3 and~5 passed $0.92$. We first assumed a bug, then an optimisation
problem, and tried longer windows, restarts with different seeds, warm-up,
cosine decay and a larger encoder. None of them helped. What helped was
replacing the 500 raw samples of a window by eleven statistics per channel per
second. With 60 training subjects, a network reading raw force traces has to
discover both what a gait cycle looks like and which of its properties separate
patients from controls; the statistics supply the first half of that problem,
and the sequence model is left to learn how the descriptors evolve. The
assignment permits either raw signals or feature-vector sequences, so this is a
legitimate design choice rather than a workaround.

The bidirectional block was our other architectural idea, and it did not
survive its own ablation. A causal scan gives the first tokens of a window
almost no context, which ought to hurt when the whole window is available at
inference. At matched encoder size, however, the causal and bidirectional
variants reached the same accuracy, and the causal one ranked subjects slightly
better. Mean-pooling may be the reason: it already aggregates over the whole
window, so the late tokens of a causal scan can carry the evidence the
classifier needs. We report this as a negative result. With 100 subjects the
difference also falls inside the bootstrap spread, so the fair conclusion is
that the experiment is inconclusive, not that the reverse direction is harmful.

The protocol comparison in Section~\ref{ssec:protocols} has a mechanical
explanation. Windows cut one second apart from the same recording are close to
duplicates of one another, and they share sensor placement, body mass and an
individual gait signature. Split them at random and most test windows have a
near neighbour in the training set, so a classifier can identify the person and
recall the label. The Random Forest makes the size of that shortcut obvious:
73\,\% when subjects are held out, 98\,\% when only windows are. The SVM-RBF
figure under the same leaky split, 88\,\%, is close to accuracies frequently
quoted for this database, which is at least suggestive about how some of them
were obtained.

That is also our reading of the gap to the literature. Our subject-wise results
sit in the 70--80\,\% band rather than above 90\,\%, and we do not think this
reflects a weaker method: the leaky column of Table~\ref{tab:protocols}
reproduces the high numbers with the same features and the same code. Work that
does enforce subject-disjoint evaluation on this database reports values
compatible with ours.

Two further caveats belong here. First, the configuration space we searched was
large -- forty-five configurations across token type, window length, encoder
size, regularisation and seed count -- and comparing them on pooled test
accuracy would make the reported number optimistic. Section~\ref{ssec:selection}
therefore selects the configuration by validation AUC alone; the difference
between the two criteria is about three accuracy points, which is worth knowing
when interpreting any single number in this report. Second, 100 subjects is
simply few: one subject is one accuracy point, and the $\pm 4$ bootstrap spread
is wide enough that most of our per-configuration differences are not
individually significant.

Beyond sample size, three limitations are worth naming. Only the normal-walk
recordings were used, so the dual-task recordings that often expose parkinsonian
gait more clearly remain unexploited. Stride detection relies on a fixed force
threshold, which is adequate for this cohort but would need revisiting for
severely impaired gait. And the stride-cycle tokens, which encode exactly the
timing variability the clinical literature emphasises, were the weakest input on
their own (66\,\%) even though the same features are complementary to the
amplitude statistics in a subject-level Random Forest (AUC $0.79$, $0.83$ and
$0.86$ for stride, amplitude and both). We take that as evidence that our fusion
of the two token families, not the features themselves, is what fell short.
